\section{Tor groups and flatness}
\label{section-tor}

\noindent
In this section we use some of the homological algebra
developed in the previous section to explain what
Tor groups are. Namely, suppose that $R$ is a ring
and that $M$, $N$ are two $R$-modules. Choose
a resolution $F_\bullet$ of $M$ by free $R$-modules.
See Lemma \ref{lemma-resolution-by-finite-free}.
Consider the homological complex
$$
F_\bullet \otimes_R N
:
\ldots
\to F_2 \otimes_R N
\to F_1 \otimes_R N
\to F_0 \otimes_R N
$$
We define $\text{Tor}^R_i(M, N)$ to be the $i$th homology
group of this complex. The following lemma explains in
what sense this is well defined.

\begin{lemma}
\label{lemma-tor-welldefined}
Let $R$ be a ring. Let $M_1, M_2, N$ be $R$-modules.
Suppose that $F_\bullet$ is a free resolution of
the module $M_1$ and that $G_\bullet$ is a free
resolution of the module $M_2$. Let $\varphi : M_1 \to M_2$
be a module map. Let $\alpha : F_\bullet \to G_\bullet$
be a map of complexes inducing $\varphi$ on
$M_1 = \Coker(d_{F, 1}) \to M_2 = \Coker(d_{G, 1})$,
see Lemma \ref{lemma-compare-resolutions}.
Then the induced maps
$$
H_i(\alpha) :
H_i(F_\bullet \otimes_R N)
\longrightarrow
H_i(G_\bullet \otimes_R N)
$$
are independent of the choice of $\alpha$. If $\varphi$
is an isomorphism, so are all the maps $H_i(\alpha)$.
If $M_1 = M_2$, $F_\bullet = G_\bullet$, and
$\varphi$ is the identity, so are all the maps $H_i(\alpha)$.
\end{lemma}

\begin{proof}
The proof of this lemma is identical to the proof of Lemma
\ref{lemma-ext-welldefined}.
\end{proof}

\noindent
Not only does this lemma imply that the Tor modules are well defined,
but it also provides for the functoriality of the constructions
$(M, N) \mapsto \text{Tor}_i^R(M, N)$ in the first variable. Of course
the functoriality in the second variable is evident. We leave it to
the reader to see that each of the $\text{Tor}_i^R$ is in fact
a functor
$$
\text{Mod}_R \times \text{Mod}_R \to \text{Mod}_R.
$$
Here $\text{Mod}_R$ denotes the category of $R$-modules, and
for the definition of the product category
see Categories, Definition \ref{categories-definition-product-category}.
Namely, given morphisms of $R$-modules $M_1 \to M_2$
and $N_1 \to N_2$ we get a commutative diagram
$$
\xymatrix{
\text{Tor}_i^R(M_1, N_1) \ar[r] \ar[d] &
\text{Tor}_i^R(M_1, N_2) \ar[d] \\
\text{Tor}_i^R(M_2, N_1) \ar[r] &
\text{Tor}_i^R(M_2, N_2) \\
}
$$

\begin{lemma}
\label{lemma-long-exact-sequence-tor}
Let $R$ be a ring and let $M$ be an $R$-module.
Suppose that $0 \to N' \to N \to N'' \to 0$ is a short
exact sequence of $R$-modules. There exists a long
exact sequence
$$
\text{Tor}_1^R(M, N')
\to \text{Tor}_1^R(M, N)
\to \text{Tor}_1^R(M, N'')
\to
M \otimes_R N'
\to M \otimes_R N
\to M \otimes_R N''
\to 0
$$
\end{lemma}

\begin{proof}
The proof of this is the same as the proof of
Lemma \ref{lemma-long-exact-seq-ext}.
\end{proof}

\noindent
Consider a homological double complex of $R$-modules
$$
\xymatrix{
\ldots \ar[r]^d &
A_{2, 0} \ar[r]^d &
A_{1, 0} \ar[r]^d &
A_{0, 0} \\
\ldots \ar[r]^d &
A_{2, 1} \ar[r]^d \ar[u]^\delta &
A_{1, 1} \ar[r]^d \ar[u]^\delta &
A_{0, 1} \ar[u]^\delta \\
\ldots \ar[r]^d &
A_{2, 2} \ar[r]^d \ar[u]^\delta &
A_{1, 2} \ar[r]^d \ar[u]^\delta &
A_{0, 2} \ar[u]^\delta \\
&
\ldots \ar[u]^\delta &
\ldots \ar[u]^\delta &
\ldots \ar[u]^\delta \\
}
$$
This means that $d_{i, j} : A_{i, j} \to A_{i-1, j}$
and $\delta_{i, j} : A_{i, j} \to A_{i, j-1}$ have the following
properties
\begin{enumerate}
\item Any composition of two $d_{i, j}$ is zero. In other words
the rows of the double complex are complexes.
\item Any composition of two $\delta_{i, j}$ is zero. In other words
the columns of the double complex are complexes.
\item For any pair $(i, j)$ we have $\delta_{i-1, j} \circ d_{i, j}
= d_{i, j-1} \circ \delta_{i, j}$. In other words, all the squares
commute.
\end{enumerate}
The correct thing to do is to associate a spectral sequence to
any such double complex. However, for the moment we can get away with
doing something slightly easier.

\medskip\noindent
Namely, for the purposes of this section only, given a double
complex $(A_{\bullet, \bullet}, d, \delta)$ set
$R(A)_j = \Coker(A_{1, j} \to A_{0, j})$ and
$U(A)_i = \Coker(A_{i, 1} \to A_{i, 0})$. (The letters
$R$ and $U$ are meant to suggest Right and Up.)
We endow $R(A)_\bullet$ with the structure of a complex
using the maps $\delta$. Similarly we endow $U(A)_\bullet$
with the structure of a complex using the maps $d$.
In other words we obtain the following huge commutative diagram
$$
\xymatrix{
\ldots \ar[r]^d &
U(A)_2 \ar[r]^d &
U(A)_1 \ar[r]^d &
U(A)_0 &
\\
\ldots \ar[r]^d &
A_{2, 0} \ar[r]^d \ar[u] &
A_{1, 0} \ar[r]^d \ar[u] &
A_{0, 0} \ar[r] \ar[u] &
R(A)_0 \\
\ldots \ar[r]^d &
A_{2, 1} \ar[r]^d \ar[u]^\delta &
A_{1, 1} \ar[r]^d \ar[u]^\delta &
A_{0, 1} \ar[r] \ar[u]^\delta &
R(A)_1 \ar[u]^\delta \\
\ldots \ar[r]^d &
A_{2, 2} \ar[r]^d \ar[u]^\delta &
A_{1, 2} \ar[r]^d \ar[u]^\delta &
A_{0, 2} \ar[r] \ar[u]^\delta &
R(A)_2 \ar[u]^\delta \\
&
\ldots \ar[u]^\delta &
\ldots \ar[u]^\delta &
\ldots \ar[u]^\delta &
\ldots \ar[u]^\delta \\
}
$$
(This is no longer a double complex of course.)
It is clear what a morphism $\Phi : (A_{\bullet, \bullet}, d, \delta)
\to (B_{\bullet, \bullet}, d, \delta)$ of double complexes
is, and it is clear that this induces morphisms of complexes
$R(\Phi) : R(A)_\bullet \to R(B)_\bullet$ and
$U(\Phi) : U(A)_\bullet \to U(B)_\bullet$.

\begin{lemma}
\label{lemma-no-spectral-sequence}
Let $(A_{\bullet, \bullet}, d, \delta)$ be a double complex such
that
\begin{enumerate}
\item Each row $A_{\bullet, j}$ is a resolution of $R(A)_j$.
\item Each column $A_{i, \bullet}$ is a resolution of $U(A)_i$.
\end{enumerate}
Then there are canonical isomorphisms
$$
H_i(R(A)_\bullet)
\cong
H_i(U(A)_\bullet).
$$
The isomorphisms are functorial with respect to morphisms
of double complexes with the properties above.
\end{lemma}

\begin{proof}
We will show that $H_i(R(A)_\bullet))$
and $H_i(U(A)_\bullet)$ are canonically
isomorphic to a third group. Namely
$$
\mathbf{H}_i(A) :=
\frac{
\{
(a_{i, 0}, a_{i-1, 1}, \ldots, a_{0, i})
\mid
d(a_{i, 0}) = \delta(a_{i-1, 1}), \ldots,
d(a_{1, i-1}) = \delta(a_{0, i})
\}}
{
\{
d(a_{i + 1, 0}) + \delta(a_{i, 1}),
d(a_{i, 1}) + \delta(a_{i-1, 2}),
\ldots,
d(a_{1, i}) + \delta(a_{0, i + 1})
\}
}
$$
Here we use the notational convention that $a_{i, j}$ denotes
an element of $A_{i, j}$. In other words, an element of $\mathbf{H}_i$
is represented by a zig-zag, represented as follows for $i = 2$
$$
\xymatrix{
a_{2, 0} \ar@{|->}[r] & d(a_{2, 0}) = \delta(a_{1, 1}) & \\
& a_{1, 1} \ar@{|->}[u] \ar@{|->}[r] & d(a_{1, 1}) = \delta(a_{0, 2}) \\
& & a_{0, 2} \ar@{|->}[u] \\
}
$$
Naturally, we divide out by ``trivial'' zig-zags, namely the submodule
generated by elements of the form $(0, \ldots, 0, -\delta(a_{t + 1, t-i}),
d(a_{t + 1, t-i}), 0, \ldots, 0)$. Note that there are canonical
homomorphisms
$$
\mathbf{H}_i(A) \to H_i(R(A)_\bullet), \quad
(a_{i, 0}, a_{i-1, 1}, \ldots, a_{0, i}) \mapsto
\text{class of image of }a_{0, i}
$$
and
$$
\mathbf{H}_i(A) \to H_i(U(A)_\bullet), \quad
(a_{i, 0}, a_{i-1, 1}, \ldots, a_{0, i}) \mapsto
\text{class of image of }a_{i, 0}
$$

\medskip\noindent
First we show that these maps are surjective.
Suppose that $\overline{r} \in H_i(R(A)_\bullet)$.
Let $r \in R(A)_i$ be a cocycle representing the
class of $\overline{r}$.
Let $a_{0, i} \in A_{0, i}$ be an element which
maps to $r$. Because $\delta(r) = 0$,
we see that $\delta(a_{0, i})$ is in the
image of $d$. Hence there exists an element
$a_{1, i-1} \in A_{1, i-1}$ such that
$d(a_{1, i-1}) = \delta(a_{0, i})$. This in turn
implies that $\delta(a_{1, i-1})$ is in the kernel
of $d$ (because $d(\delta(a_{1, i-1})) = \delta(d(a_{1, i-1}))
= \delta(\delta(a_{0, i})) = 0$. By exactness of the
rows we find an element $a_{2, i-2}$ such that
$d(a_{2, i-2}) = \delta(a_{1, i-1})$. And so on
until a full zig-zag is found. Of course surjectivity
of $\mathbf{H}_i \to H_i(U(A))$ is shown similarly.

\medskip\noindent
To prove injectivity we argue in exactly the same way.
Namely, suppose we are given a zig-zag
$(a_{i, 0}, a_{i-1, 1}, \ldots, a_{0, i})$
which maps to zero in $H_i(R(A)_\bullet)$.
This means that $a_{0, i}$ maps to an element
of $\Coker(A_{i, 1} \to A_{i, 0})$
which is in the image of
$\delta : \Coker(A_{i + 1, 1} \to A_{i + 1, 0}) \to
\Coker(A_{i, 1} \to A_{i, 0})$.
In other words, $a_{0, i}$ is in the image of
$\delta \oplus d : A_{0, i + 1} \oplus A_{1, i} \to A_{0, i}$.
From the definition of trivial zig-zags we see that
we may modify our zig-zag by a trivial one and
assume that $a_{0, i} = 0$. This immediately
implies that $d(a_{1, i-1}) = 0$. As the rows
are exact this implies that $a_{1, i-1}$ is
in the image of $d : A_{2, i-1} \to A_{1, i-1}$.
Thus we may modify our zig-zag once again by a
trivial zig-zag and assume that our zig-zag looks
like $(a_{i, 0}, a_{i-1, 1}, \ldots, a_{2, i-2}, 0, 0)$.
Continuing like this we obtain the desired injectivity.

\medskip\noindent
If $\Phi : (A_{\bullet, \bullet}, d, \delta)
\to (B_{\bullet, \bullet}, d, \delta)$ is a morphism
of double complexes both of which satisfy the conditions
of the lemma, then we clearly obtain a commutative
diagram
$$
\xymatrix{
H_i(U(A)_\bullet) \ar[d] &
\mathbf{H}_i(A) \ar[r] \ar[l] \ar[d] &
H_i(R(A)_\bullet) \ar[d] \\
H_i(U(B)_\bullet) &
\mathbf{H}_i(B) \ar[r] \ar[l] &
H_i(R(B)_\bullet) \\
}
$$
This proves the functoriality.
\end{proof}

\begin{remark}
\label{remark-signs-double-complex}
The isomorphism constructed above is the ``correct'' one only up to signs.
A good part of homological algebra is concerned with choosing signs for
various maps and showing commutativity of diagrams with intervention
of suitable signs. For the moment we will simply use the isomorphism
as given in the proof above, and worry about signs later.
\end{remark}

\begin{lemma}
\label{lemma-tor-left-right}
Let $R$ be a ring. For any $i \geq 0$ the functors
$\text{Mod}_R \times \text{Mod}_R \to \text{Mod}_R$,
$(M, N) \mapsto \text{Tor}_i^R(M, N)$ and
$(M, N) \mapsto \text{Tor}_i^R(N, M)$ are
canonically isomorphic.
\end{lemma}

\begin{proof}
Let $F_\bullet$ be a free resolution of the module $M$ and
let $G_\bullet$ be a free resolution of the module $N$.
Consider the double complex $(A_{i, j}, d, \delta)$ defined
as follows:
\begin{enumerate}
\item set $A_{i, j} = F_i \otimes_R G_j$,
\item set $d_{i, j} : F_i \otimes_R G_j \to F_{i-1} \otimes G_j$
equal to $d_{F, i} \otimes \text{id}$, and
\item set $\delta_{i, j} : F_i \otimes_R G_j \to F_i \otimes G_{j-1}$
equal to $\text{id} \otimes d_{G, j}$.
\end{enumerate}
This double complex is usually simply denoted $F_\bullet \otimes_R G_\bullet$.

\medskip\noindent
Since each $G_j$ is free, and hence flat we see that each
row of the double complex is exact except in homological
degree $0$. Since each $F_i$ is free and hence flat we see that each
column of the double complex is exact except in homological
degree $0$. Hence the double complex satisfies the conditions
of Lemma \ref{lemma-no-spectral-sequence}.

\medskip\noindent
To see what the lemma says we compute $R(A)_\bullet$ and $U(A)_\bullet$.
Namely,
\begin{eqnarray*}
R(A)_i & = & \Coker(A_{1, i} \to A_{0, i}) \\
& = & \Coker(F_1 \otimes_R G_i \to F_0 \otimes_R G_i) \\
& = & \Coker(F_1 \to F_0) \otimes_R G_i \\
& = & M \otimes_R G_i
\end{eqnarray*}
In fact these isomorphisms are compatible with the differentials
$\delta$ and we see that $R(A)_\bullet = M \otimes_R G_\bullet$
as homological complexes. In exactly the same way we see that
$U(A)_\bullet = F_\bullet \otimes_R N$. We get
\begin{eqnarray*}
\text{Tor}_i^R(M, N)
& = & H_i(F_\bullet \otimes_R N) \\
& = & H_i(U(A)_\bullet) \\
& = & H_i(R(A)_\bullet) \\
& = & H_i(M \otimes_R G_\bullet) \\
& = & H_i(G_\bullet \otimes_R M) \\
& = & \text{Tor}_i^R(N, M)
\end{eqnarray*}
Here the third equality is Lemma \ref{lemma-no-spectral-sequence}, and
the fifth equality uses the isomorphism $V \otimes W = W \otimes V$
of the tensor product.

\medskip\noindent
Functoriality. Suppose that we have $R$-modules $M_\nu$, $N_\nu$,
$\nu = 1, 2$. Let $\varphi : M_1 \to M_2$ and $\psi : N_1 \to N_2$
be morphisms of $R$-modules.
Suppose that we have free resolutions $F_{\nu, \bullet}$
for $M_\nu$ and free resolutions $G_{\nu, \bullet}$ for $N_\nu$.
By Lemma \ref{lemma-compare-resolutions} we may choose
maps of complexes $\alpha : F_{1, \bullet} \to F_{2, \bullet}$
and $\beta : G_{1, \bullet} \to G_{2, \bullet}$ compatible
with $\varphi$ and $\psi$. We claim that
the pair $(\alpha, \beta)$ induces a morphism of double
complexes
$$
\alpha \otimes \beta :
F_{1, \bullet} \otimes_R G_{1, \bullet}
\longrightarrow
F_{2, \bullet} \otimes_R G_{2, \bullet}
$$
This is really a very straightforward check using the rule
that $F_{1, i} \otimes_R G_{1, j} \to F_{2, i} \otimes_R G_{2, j}$
is given by $\alpha_i \otimes \beta_j$ where $\alpha_i$,
resp.\ $\beta_j$ is the degree $i$, resp.\ $j$ component of $\alpha$,
resp.\ $\beta$. The reader also readily verifies that the
induced maps $R(F_{1, \bullet} \otimes_R G_{1, \bullet})_\bullet
\to R(F_{2, \bullet} \otimes_R G_{2, \bullet})_\bullet$
agrees with the map $M_1 \otimes_R G_{1, \bullet}
\to M_2 \otimes_R G_{2, \bullet}$ induced by $\varphi \otimes \beta$.
Similarly for the map induced on the $U(-)_\bullet$ complexes.
Thus the statement on functoriality follows from the statement
on functoriality in Lemma \ref{lemma-no-spectral-sequence}.
\end{proof}

\begin{remark}
\label{remark-curiosity-signs-swap}
An interesting case occurs when $M = N$ in the above.
In this case we get a canonical map $\text{Tor}_i^R(M, M)
\to \text{Tor}_i^R(M, M)$. Note that this map is not the
identity, because even when $i = 0$ this map is not the
identity! For example, if $V$ is a vector space of dimension
$n$ over a field, then the switch map $V \otimes_k V \to V \otimes_k V$
has $(n^2 + n)/2$ eigenvalues $+1$ and $(n^2-n)/2$ eigenvalues
$-1$. In characteristic $2$ it is not even diagonalizable.
Note that even changing the sign of the map will not get rid
of this.
\end{remark}

\begin{lemma}
\label{lemma-tor-noetherian}
Let $R$ be a Noetherian ring. Let $M$, $N$ be finite $R$-modules.
Then $\text{Tor}_p^R(M, N)$ is a finite $R$-module for all $p$.
\end{lemma}

\begin{proof}
This holds because $\text{Tor}_p^R(M, N)$ is computed as the
cohomology groups of a complex $F_\bullet \otimes_R N$
with each $F_n$ a finite free $R$-module, see
Lemma \ref{lemma-resolution-by-finite-free}.
\end{proof}

\begin{lemma}
\label{lemma-characterize-flat}
Let $R$ be a ring. Let $M$ be an $R$-module.
The following are equivalent:
\begin{enumerate}
\item The module $M$ is flat over $R$.
\item For all $i > 0$ the functor $\text{Tor}_i^R(M, -)$ is zero.
\item The functor $\text{Tor}_1^R(M, -)$ is zero.
\item For all ideals $I \subset R$ we have $\text{Tor}_1^R(M, R/I) = 0$.
\item For all finitely generated ideals $I \subset R$ we have
$\text{Tor}_1^R(M, R/I) = 0$.
\end{enumerate}
\end{lemma}

\begin{proof}
Suppose $M$ is flat. Let $N$ be an $R$-module.
Let $F_\bullet$ be a free resolution of $N$.
Then $F_\bullet \otimes_R M$ is a resolution of $N \otimes_R M$,
by flatness of $M$. Hence all higher Tor groups vanish.

\medskip\noindent
It now suffices to show that the last condition implies that
$M$ is flat. Let $I \subset R$ be an ideal.
Consider the short exact sequence
$0 \to I \to R \to R/I \to 0$. Apply
Lemma \ref{lemma-long-exact-sequence-tor}. We get an
exact sequence
$$
\text{Tor}_1^R(M, R/I) \to
M \otimes_R I \to
M \otimes_R R \to
M \otimes_R R/I \to
0
$$
Since obviously $M \otimes_R R = M$ we conclude that the
last hypothesis implies that $M \otimes_R I \to M$ is
injective for every finitely generated ideal $I$.
Thus $M$ is flat by Lemma \ref{lemma-flat}.
\end{proof}

\begin{remark}
\label{remark-Tor-ring-mod-ideal}
The proof of Lemma \ref{lemma-characterize-flat} actually shows
that
$$
\text{Tor}_1^R(M, R/I)
=
\Ker(I \otimes_R M \to M).
$$
\end{remark}














